\section{Verilog codes}
\setcounter{lstlisting}{0}
\renewcommand{\thelstlisting}{\thesection.\arabic{lstlisting}}
\label{appendix:verilog_codes}

\definecolor{verilogcommentcolor}{RGB}{104,180,104}
\definecolor{verilogkeywordcolor}{RGB}{49,49,255}
\definecolor{verilogsystemcolor}{RGB}{128,0,255}
\definecolor{verilognumbercolor}{RGB}{255,143,102}
\definecolor{verilogstringcolor}{RGB}{160,160,160}
\definecolor{verilogdefinecolor}{RGB}{128,64,0}
\definecolor{verilogoperatorcolor}{RGB}{0,0,128}

% Verilog style
\lstdefinestyle{prettyverilog}{
   language           = Verilog,
   commentstyle       = \color{verilogcommentcolor},
   alsoletter         = \$'0123456789\`,
   literate           = *{+}{{\verilogColorOperator{+}}}{1}%
                         {-}{{\verilogColorOperator{-}}}{1}%
                         {@}{{\verilogColorOperator{@}}}{1}%
                         {;}{{\verilogColorOperator{;}}}{1}%
                         {*}{{\verilogColorOperator{*}}}{1}%
                         {?}{{\verilogColorOperator{?}}}{1}%
                         {:}{{\verilogColorOperator{:}}}{1}%
                         {<}{{\verilogColorOperator{<}}}{1}%
                         {>}{{\verilogColorOperator{>}}}{1}%
                         {=}{{\verilogColorOperator{=}}}{1}%
                         {!}{{\verilogColorOperator{!}}}{1}%
                         {^}{{\verilogColorOperator{$\land$}}}{1}%
                         {|}{{\verilogColorOperator{|}}}{1}%
                         {=}{{\verilogColorOperator{=}}}{1}%
                         {[}{{\verilogColorOperator{[}}}{1}%
                         {]}{{\verilogColorOperator{]}}}{1}%
                         {(}{{\verilogColorOperator{(}}}{1}%
                         {)}{{\verilogColorOperator{)}}}{1}%
                         {,}{{\verilogColorOperator{,}}}{1}%
                         {.}{{\verilogColorOperator{.}}}{1}%
                         {~}{{\verilogColorOperator{$\sim$}}}{1}%
                         {\%}{{\verilogColorOperator{\%}}}{1}%
                         {\&}{{\verilogColorOperator{\&}}}{1}%
                         {\#}{{\verilogColorOperator{\#}}}{1}%
                         {\ /\ }{{\verilogColorOperator{\ /\ }}}{3}%
                         {\ _}{\ \_}{2}%
                        ,
   morestring         = [s][\color{verilogstringcolor}]{"}{"},%
   identifierstyle    = \color{black},
   vlogdefinestyle    = \color{verilogdefinecolor},
   vlogconstantstyle  = \color{verilognumbercolor},
   vlogsystemstyle    = \color{verilogsystemcolor},
   basicstyle         = \scriptsize\fontencoding{T1}\ttfamily,
   keywordstyle       = \bfseries\color{verilogkeywordcolor},
   numbers            = left,
   numbersep          = 10pt,
   tabsize            = 4,
   escapeinside       = {/*!}{!*/},
   upquote            = true,
   sensitive          = true,
   showstringspaces   = false, %without this there will be a symbol in the places where there is a space
   frame              = single
}


% This is shamelessly stolen and modified from:
% https://github.com/jubobs/sclang-prettifier/blob/master/sclang-prettifier.dtx
\makeatletter

% Language name
\newcommand\language@verilog{Verilog}
\expandafter\lst@NormedDef\expandafter\languageNormedDefd@verilog%
  \expandafter{\language@verilog}
  
% save definition of single quote for testing
\lst@SaveOutputDef{`'}\quotesngl@verilog
\lst@SaveOutputDef{``}\backtick@verilog
\lst@SaveOutputDef{`\$}\dollar@verilog

% Extract first character token in sequence and store in macro 
% firstchar@verilog, per http://tex.stackexchange.com/a/159267/21891
\newcommand\getfirstchar@verilog{}
\newcommand\getfirstchar@@verilog{}
\newcommand\firstchar@verilog{}
\def\getfirstchar@verilog#1{\getfirstchar@@verilog#1\relax}
\def\getfirstchar@@verilog#1#2\relax{\def\firstchar@verilog{#1}}

% Initially empty hook for lst
\newcommand\addedToOutput@verilog{}
\lst@AddToHook{Output}{\addedToOutput@verilog}

% The style used for constants as set in lstdefinestyle
\newcommand\constantstyle@verilog{}
\lst@Key{vlogconstantstyle}\relax%
   {\def\constantstyle@verilog{#1}}

% The style used for defines as set in lstdefinestyle
\newcommand\definestyle@verilog{}
\lst@Key{vlogdefinestyle}\relax%
   {\def\definestyle@verilog{#1}}

% The style used for defines as set in lstdefinestyle
\newcommand\systemstyle@verilog{}
\lst@Key{vlogsystemstyle}\relax%
   {\def\systemstyle@verilog{#1}}

% Counter used to check current character is a digit
\newcount\currentchar@verilog
  
% Processing macro
\newcommand\@ddedToOutput@verilog
{%
   % If we're in \lstpkg{}' processing mode...
   \ifnum\lst@mode=\lst@Pmode%
      % Save the first token in the current identifier to \@getfirstchar
      \expandafter\getfirstchar@verilog\expandafter{\the\lst@token}%
      % Check if the token is a backtick
      \expandafter\ifx\firstchar@verilog\backtick@verilog
         % If so, then this starts a define
         \let\lst@thestyle\definestyle@verilog%
      \else
         % Check if the token is a dollar
         \expandafter\ifx\firstchar@verilog\dollar@verilog
            % If so, then this starts a system command
            \let\lst@thestyle\systemstyle@verilog%
         \else
            % Check if the token starts with a single quote
            \expandafter\ifx\firstchar@verilog\quotesngl@verilog
               % If so, then this starts a constant without length
               \let\lst@thestyle\constantstyle@verilog%
            \else
               \currentchar@verilog=48
               \loop
                  \expandafter\ifnum%
                  \expandafter`\firstchar@verilog=\currentchar@verilog%
                     \let\lst@thestyle\constantstyle@verilog%
                     \let\iterate\relax%
                  \fi
                  \advance\currentchar@verilog by \@ne%
                  \unless\ifnum\currentchar@verilog>57%
               \repeat%
            \fi
         \fi
      \fi
      % ...but override by keyword style if a keyword is detected!
      %\lsthk@DetectKeywords% 
   \fi
}

% Add processing macro only if verilog
\lst@AddToHook{PreInit}{%
  \ifx\lst@language\languageNormedDefd@verilog%
    \let\addedToOutput@verilog\@ddedToOutput@verilog%
  \fi
}

% Colour operators in literate
\newcommand{\verilogColorOperator}[1]
{%
  \ifnum\lst@mode=\lst@Pmode\relax%
   {\bfseries\textcolor{verilogoperatorcolor}{#1}}%
  \else
    #1%
  \fi
}

\makeatother
%%%%%%%%%%%%%%%%%%%%%%%%%%%%%%%%%%%%%%%%%%%%%%%%%%%%%%%%%%%%%%%%%%%%%%%%%%%%%%%%%%%%%%%%%%%%%%%%%%%%%%%%%%%%%%%%%%%%%%%%%%%%%
% End Verilog Code Style
%%%%%%%%%%%%%%%%%%%%%%%%%%%%%%%%%%%%%%%%%%%%%%%%%%%%%%%%%%%%%%%%%%%%%%%%%%%%%%%%%%%%%%%%%%%%%%%%%%%%%%%%%%%%%%%%%%%%%%%%%%%%%


\begin{lstlisting}[style={prettyverilog}, caption={Verilog code for the 4-bit convergent quantizer with gain and dynamic shift.}, label={lst:verilog_quantizer}]
module quantizer (
    input  wire          clk,
    input  wire          ce,
    input  wire signed [30:0] din,   
    input  wire [15:0]   gain,       
    input  wire [5:0]    shift,      
    output reg signed [3:0] dout,
    output reg           clip_flag
);
    // Stage 1: multiplication
    reg signed [47:0] mult_result;
    // Stage 2: shift left
    reg signed [47:0] shifted_stage;
    // Stage 3: rounding variables
    reg signed [47:0] rounded_stage; 
    reg pattern_detect;
    reg signed [47:0] multadd_reg;
    
    // Delays for shifted_stage to align with rounding pipeline
    reg signed [47:0] shifted_stage_d1;  
    reg signed [47:0] shifted_stage_d2;  
    
    // Convergent rounding constants
    wire [43:0] pattern = 44'b00000000000000000000000000000000000000000000;
    wire [47:0] c = 48'b000001111111111111111111111111111111111111111111; 
    wire signed [47:0] multadd;
    
    // Hierarchical OR for fractional part detection (hardware efficient)
    wire frac_nonzero_1 = |shifted_stage_d2[43:33];  // 11 bits
    wire frac_nonzero_2 = |shifted_stage_d2[32:22];  // 11 bits  
    wire frac_nonzero_3 = |shifted_stage_d2[21:11];  // 11 bits
    wire frac_nonzero_4 = |shifted_stage_d2[10:0];   // 11 bits
    wire any_frac_bit = frac_nonzero_1 | frac_nonzero_2 | frac_nonzero_3 | frac_nonzero_4;
    
    // =========================
    // Stage 1: Multiply by gain
    // =========================
    always @(posedge clk) begin
        if (ce)
            mult_result <= din * $signed({1'b0, gain});
        else
            mult_result <= 0;
    end
    
    // =========================
    // Stage 2: Shift LEFT
    // =========================
    always @(posedge clk) begin
        if (ce)
            shifted_stage <= mult_result <<< shift;  // shift left
        else
            shifted_stage <= 0;
    end
    
    // =========================
    // Delay registers for shifted_stage
    // =========================
    always @(posedge clk) begin
        if (ce) begin
            shifted_stage_d1 <= shifted_stage;    
            shifted_stage_d2 <= shifted_stage_d1; 
        end else begin
            shifted_stage_d1 <= 0;
            shifted_stage_d2 <= 0;
        end
    end
    
    // =========================
    // Stage 3: Convergent Rounding Logic
    // =========================
    assign multadd = shifted_stage + c + 1'b1;
    
    always @(posedge clk) begin
        if (ce) begin
            pattern_detect <= (multadd[43:0] == pattern) ? 1'b1 : 1'b0;
            multadd_reg <= multadd;
        end else begin
            pattern_detect <= 1'b0;
            multadd_reg <= 48'b0;
        end
    end
    
    always @(posedge clk) begin
        if (ce) begin
            // Convergent rounding: if pattern detected (midpoint), force LSB to 0
            if (pattern_detect)
                rounded_stage <= {multadd_reg[47:45], 1'b0, multadd_reg[43:0]}; 
            else
                rounded_stage <= multadd_reg; // Normal half up rounding
        end else begin
            rounded_stage <= 48'b0;
        end
    end
    
    // =========================
    // Stage 4: Overflow Detection and Saturation
    // =========================
    always@(posedge clk) begin
        if (ce) begin
            // Use delayed shifted_stage aligned with rounded_stage
            if ((shifted_stage_d2[47:44] == 4'b0111) && any_frac_bit) begin
                // Positive overflow: shifted_stage > 7.0
                dout <= 4'b0111; // Max positive value for 4-bit signed
                clip_flag <= 1'b1;
            end else begin
                // No overflow, assign the rounded value
                dout <= rounded_stage[47:44]; // Take the top 4 bits as output
                clip_flag <= 1'b0;
            end
        end else begin
            dout <= 4'b0;
            clip_flag <= 1'b0;
        end
    end
    
endmodule
\end{lstlisting}

\begin{lstlisting}[style={prettyverilog}, caption={Verilog code for FFT reorder module.}, label={lst:verilog_fft_reorder}]
module fft_reorderer (
    input wire clk,
    input wire ce,          
    input wire rst,         
    
    // Control
    input wire sync,        // Start pulse (edge detected)
    output reg sync_out = 0,// Start of Stream pulse for next module
    
    // Data
    input wire [63:0] din,
    output reg [63:0] dout = 0,
    output reg dout_valid = 0
);

    localparam WORDS_PER_FRAME = 1024;
    localparam BRAM_DEPTH      = 2048; // 2 Frames for Ping-Pong

    reg [9:0] cnt              = 0;
    reg       active           = 0;
    reg       ping_pong        = 0; // 0: Wr A / Rd B,  1: Wr B / Rd A
    reg       first_frame_done = 0;

    // Edge detector for SYNC
    reg sync_d = 0;
    wire sync_pulse;

    // Control pipeline (latency compensation)
    reg read_en_d1 = 0;

    // BRAM Interface
    reg [10:0]  bram_addr_wr = 0;
    reg [10:0]  bram_addr_rd = 0;
    wire [63:0] bram_dout_raw;

    assign sync_pulse = sync && !sync_d;

    // ============================================================
    // DUAL PORT BRAM INSTANCE
    // ============================================================
    bram_dual_port #(
        .DEPTH(BRAM_DEPTH),
        .WIDTH(64)
    ) my_bram (
        .clk   (clk),
        .ena   (ce),
        .enb   (ce),
        .wea   (active),
        .addra (bram_addr_wr),
        .addrb (bram_addr_rd),
        .dia   (din),
        .dob   (bram_dout_raw)
    );

    // ============================================================
    // MAIN LOGIC
    // ============================================================
    always @(posedge clk) begin
        if (rst) begin
            cnt              <= 0;
            active           <= 0;
            ping_pong        <= 0;
            first_frame_done <= 0;
            sync_d           <= 0;
            read_en_d1       <= 0;
            dout             <= 0;
            dout_valid       <= 0;
            sync_out         <= 0;
            bram_addr_wr     <= 0;
            bram_addr_rd     <= 0;
            
        end else if (ce) begin
            sync_d <= sync;

            // Start management
            if (sync_pulse) begin
                active           <= 1;
                cnt              <= 0;
                ping_pong        <= 0;
                first_frame_done <= 0;
                read_en_d1       <= 0;
                dout_valid       <= 0;
                sync_out         <= 0;
            end
            
            else if (active) begin
                
                // Address calculation
                // Write: Base (0 or 1024) + Counter
                bram_addr_wr <= (ping_pong ? WORDS_PER_FRAME : 0) + cnt;

                // Read with FFT Shift: Opposite base + Counter with MSB inverted
                // Inverting cnt[9] swaps lower and upper halves (0-511 <-> 512-1023)
                bram_addr_rd <= (ping_pong ? 0 : WORDS_PER_FRAME) + { ~cnt[9], cnt[8:0] };

                // State counters
                if (cnt == WORDS_PER_FRAME - 1) begin
                    cnt              <= 0;
                    ping_pong        <= ~ping_pong;
                    first_frame_done <= 1;
                end else begin
                    cnt              <= cnt + 1;
                end

                // Output pipeline (total latency = 2 cycles)
                read_en_d1 <= first_frame_done;

                // sync_out: Assert 1 cycle before first valid data output
                sync_out <= (first_frame_done && !read_en_d1);

                dout       <= bram_dout_raw;
                dout_valid <= read_en_d1;
            end 
            else begin
                dout_valid <= 0;
                sync_out   <= 0;
            end
        end
    end

endmodule

\end{lstlisting}

\begin{lstlisting}[style={prettyverilog}, caption={Verilog code for dual-port BRAM used in \texttt{fft\_reorderer.v}.}, label={lst:verilog_dual_port_bram}]
// Simple Dual-Port Block RAM with One Clock UG901
module bram_dual_port #(
    parameter DEPTH = 2048,
    parameter WIDTH = 64,
    parameter ADDR_WIDTH = $clog2(DEPTH)
)(
    input wire clk,
    input wire ena,
    input wire enb,
    input wire wea,
    input wire [ADDR_WIDTH-1:0] addra,
    input wire [ADDR_WIDTH-1:0] addrb,
    input wire [WIDTH-1:0] dia,
    output reg [WIDTH-1:0] dob
);

    (* ram_style = "block" *) reg [WIDTH-1:0] ram [0:DEPTH-1];

    // Port A: Write
    always @(posedge clk) begin
        if (ena && wea)
            ram[addra] <= dia;
    end

    // Port B: Read
    always @(posedge clk) begin
        if (enb)
            dob <= ram[addrb];
    end

endmodule
\end{lstlisting}

\begin{lstlisting}[style={prettyverilog}, caption={Verilog code for the LIFO module.}, label={lst:verilog_lifo}]
module lifo (
    input wire clk,
    input wire ce,
    input wire [63:0] din,
    input wire we,
    input wire re,
    input wire rst,
    output reg [63:0] dout,
    output reg empty,
    output reg full
);
    // Parameters
    parameter integer DATA_WIDTH = 64;
    parameter integer DEPTH = 256;

    // Function to calculate log2 (ceiling)
    function integer clog2;
        input integer value;
        begin
            value = value - 1;
            for (clog2 = 0; value > 0; clog2 = 0 + 1)
                value = value >> 1;
        end
    endfunction

    // Pointer width: 8 bits for address + 1 extra bit to distinguish full/empty
    localparam integer PTR_WIDTH = clog2(DEPTH) + 1;

    // Block RAM memory
    (* ram_style = "block" *) reg [DATA_WIDTH - 1:0] mem [0:DEPTH - 1];
    
    // Stack pointer
    reg [PTR_WIDTH - 1:0] ptr;

    // Status flags
    always @(*) begin
        full  = (ptr == DEPTH);
        empty = (ptr == 0);
    end

    // BRAM write and read operations
    always @(posedge clk) begin
        // Push: Write data to current pointer position
        if (we && !full && !rst) begin
            mem[ptr] <= din;
        end
        
        // Pop: Read from top of stack (ptr-1)
        if (re && !empty && !rst) begin
            dout <= mem[ptr - 1]; 
        end
    end

    // Pointer control
    initial begin
        ptr = 0;
        dout = 0;
    end

    always @(posedge clk) begin
        if (rst) begin
            ptr <= 0;
        end
        else begin
            if (we && !full) begin
                ptr <= ptr + 1;
            end
            else if (re && !empty) begin
                ptr <= ptr - 1;
            end
        end
    end

endmodule
    
\end{lstlisting}


\begin{lstlisting}[style={prettyverilog}, caption={Verilog code for the skid buffer module used in \texttt{piso.v}.}, label={lst:verilog_skid_buffer}]
module skid_buffer #(
    parameter DIN_WIDTH = 64
) (
    input wire clk,
    input wire rst,

    input wire [DIN_WIDTH-1:0] din,
    input wire din_valid, 
    output wire  din_ready, 

    output wire dout_valid, 
    input wire dout_ready, 
    output wire [DIN_WIDTH-1:0] dout
);

reg [DIN_WIDTH-1:0] din_r=0;

//register valid
reg val=0;
always@(posedge clk)begin
    if(rst)
        val <=0;
    else if((din_valid & din_ready) && (dout_valid & ~dout_ready)) begin
        //there is din valid but dout is stalled
        val <=1;
    end 
    else if(dout_ready)
        val <= 0;
end


//register data
always@(posedge clk)begin
    if(rst)
        din_r <=0;
    else if(din_valid & din_ready)
        din_r <= din;
end

//din ready
assign din_ready = ~val;

//dout side
reg [DIN_WIDTH-1:0] dout_r=0;
reg dout_valid_r=0;


//valid
reg flag =0;
wire flag2 = din_valid | val;
always@(posedge clk)begin
    if(rst)
        dout_valid_r <=0;
    else if((~dout_valid) | dout_ready)begin
        dout_valid_r <= (din_valid | val);
        flag <= 1;
    end
    else
        flag <= 0;
end

//data
always@(posedge clk)begin
    if(rst)
        dout_r <= 0;
    else if((~dout_valid) | dout_ready)begin
        if(val)
            dout_r <= din_r;
        else if(din_valid)
            dout_r <= din;
        else
            dout_r <= 0;
    end
end

assign dout = dout_r;
assign dout_valid = dout_valid_r; 

endmodule
\end{lstlisting}

\begin{lstlisting}[style={prettyverilog}, caption={Verilog code for PISO converter used in \texttt{hge\_write\_packetizer.v}.}, label={lst:verilog_sipo_converter}]
/* DIN --> FIFO --> time multiplex -->DOUT */

module piso #(
    parameter DIN_WIDTH = 512,
    parameter DOUT_WIDTH = 64,
    parameter FIFO_DEPTH = 64,
    //stupid ise dont allow $clog2 in localparam
    parameter CYCLES = DIN_WIDTH/DOUT_WIDTH,
    parameter TIME_SIZE = $clog2(CYCLES)
) (
    input wire clk,
    input wire rst,

    input wire [DIN_WIDTH-1:0] din,
    input wire din_valid,

    output wire [DOUT_WIDTH-1:0] dout,
    output wire dout_valid,
    input wire dout_ready,

    output wire fifo_full
);

//localparam CYCLES = DIN_WIDTH/DOUT_WIDTH;

/***
 ***   FIFO
 ***/
reg [$clog2(FIFO_DEPTH):0] waddr=0, raddr=0;
wire ren;
wire [DIN_WIDTH-1:0] dout_fifo;

//write pointer
always@(posedge clk)begin
    if(rst)
        waddr <=0;
    else if(din_valid & ~full)
        waddr <= waddr+1;
end

bram_infer #(
    .N_ADDR(FIFO_DEPTH),
    .DATA_WIDTH(DIN_WIDTH)
) bram_inst (
    .clk(clk),
    .wen(din_valid),
    .ren(ren),
    .wadd(waddr[$clog2(FIFO_DEPTH)-1:0]),
    .radd(raddr[$clog2(FIFO_DEPTH)-1:0]),
    .win(din),
    .wout(dout_fifo)
);
//fifo control signals

wire empty, full;
assign empty = (waddr==raddr);
assign full = ((waddr[$clog2(FIFO_DEPTH)] != (raddr[$clog2(FIFO_DEPTH)])) &
            (waddr[$clog2(FIFO_DEPTH)-1:0]== raddr[$clog2(FIFO_DEPTH)-1:0]));

//read pointer
always@(posedge clk)begin
    if(rst)
        raddr <=0;
    else if(ren)
        raddr <= raddr+1;
end

assign ren = (~empty & (state==IDLE) & sk_ready);

/***
 ***    Time multiplex
 ***/

reg [$clog2(CYCLES):0] time_multiplex=0;
localparam IDLE = 1'b0;
localparam BUSY = 1'b1;
reg state=0, next_state=0;

always@(posedge clk)begin
    if(rst)
        state <= IDLE;
    else
        state <= next_state;
end

always@(*)begin
    case(state)
        IDLE: begin
            if(~empty & sk_ready)  next_state = BUSY;
            else        next_state = IDLE;
        end
        BUSY: begin
            if((time_multiplex==(CYCLES-1)) & ~stall)   next_state = IDLE; 
            else                                        next_state = BUSY;
        end
    endcase
end

//control signals
wire stall = sk_valid & ~sk_ready;
reg dout_valid_r=0;
always@(posedge clk)begin
    case(state)
        IDLE: begin
            //time_multiplex <={($clog2(CYCLES)+1){1'b1}};
            time_multiplex <={(TIME_SIZE+1){1'b1}};
            dout_valid_r <=0;
        end
        BUSY:begin
            if(stall)
                dout_valid_r <=1;
            else begin 
                time_multiplex <= time_multiplex+1;
                dout_valid_r <=1;
            end
        end
    endcase
end

reg [DOUT_WIDTH-1:0] serial_out = 0;
always@(*)
    serial_out = dout_fifo[DOUT_WIDTH*time_multiplex+:DOUT_WIDTH];

/*
reg dout_valid_rr=0;
always@(posedge clk)begin
    serial_out <= dout_fifo[DOUT_WIDTH*time_multiplex+:DOUT_WIDTH];
    dout_valid_rr <= dout_valid_r;
end
*/
wire sk_ready, sk_valid;
assign sk_valid = dout_valid_r & (state==BUSY);

skid_buffer #(
    .DIN_WIDTH(DOUT_WIDTH)
)skid_buffer_inst (
    .clk(clk),
    .rst(rst),
    .din(serial_out),
    .din_valid(sk_valid), 
    .din_ready(sk_ready), 
    .dout_valid(dout_valid), 
    .dout_ready(dout_ready), 
    .dout(dout)
);

assign fifo_full = full;
endmodule
\end{lstlisting}

\begin{lstlisting}[style={prettyverilog}, caption={Verilog code for the HGE write packetizer used in \texttt{packetizer\_top.v}.}, label={lst:verilog_hge_packetizer}]
module hge_write_packetizer #(
    parameter DIN_WIDTH = 2048,
    parameter FIFO_DEPTH = 512
) (
    input wire clk,
    input wire rst,
    input wire [DIN_WIDTH-1:0] din,
    input wire din_valid,
    
    // Configuration signals and timestamp
    input wire [31:0] pkt_len,      
    input wire [31:0] sleep_cycles,
    input wire [31:0] timestamp_sec,  
    input wire [31:0] timestamp_usec, 

    // To the 100Gbe block 
    output wire [511:0] tx_data,
    output wire tx_valid,
    output wire tx_eof,

    output wire fifo_full
);

    wire [511:0] tge_data;
    wire tge_data_valid;

    // Definition of states
    localparam WAIT   = 0;
    localparam HEADER = 1; 
    localparam READ   = 2;
    
    reg [1:0] state = WAIT;
    reg [1:0] next_state;

    reg [31:0] counter = 0;
    
    // Counters for the Header
    reg [63:0] seq_num = 0;  
    reg [1:0]  band_idx = 0; 

    // Logic for the Next State
    always@(*) begin
        case(state)
            WAIT: begin
                if(counter == sleep_cycles) next_state = HEADER;
                else                        next_state = WAIT;
            end
            HEADER: begin
                // Inject 1 cycle of header and move to data
                next_state = READ; 
            end
            READ: begin
                if(counter == pkt_len)      next_state = WAIT;
                else                        next_state = READ;
            end
            default: next_state = WAIT;
        endcase
    end

    // Control and output signals
    wire valid_req = tge_data_valid & piso_ready;
    reg eof = 0;
    reg [511:0] dout;
    reg dout_valid = 0;

    // Header structure:
    // [Padding (382b)] [Sec (32b)] [uSec (32b)] [BandIdx (2b)] [SeqNum (64b)]
    // Total bits used: 32 + 32 + 8 + 64 = 136 bits.
    // Padding: 512 - 136 = 376 bits.
    
    wire [511:0] header_word;
    assign header_word = {376'd0, timestamp_sec, timestamp_usec, 6'd0, band_idx, seq_num};

    always @(posedge clk) begin
        if (rst) begin
            counter    <= 0;
            state      <= WAIT;
            eof        <= 0;
            dout_valid <= 0;
            seq_num    <= 0;
            band_idx   <= 0;
        end
        else begin
            state <= next_state; 
            
            case(state)
                WAIT: begin
                    eof        <= 0;
                    dout_valid <= 0;
                    if(counter == sleep_cycles) begin
                        counter <= 0;
                    end
                    else begin
                        counter <= counter + 1;
                    end
                end

                HEADER: begin
                    // 1. Send Header (capture current timestamps)
                    dout       <= header_word;
                    dout_valid <= 1;
                    eof        <= 0; 
                    
                    counter    <= 0; 
                end

                READ: begin
                    dout <= tge_data;
                    
                    if (valid_req) begin
                        dout_valid <= 1;
                        if (counter == pkt_len) begin
                            eof     <= 1;
                            counter <= 0;
                            
                            // Update Header Counters
                            seq_num <= seq_num + 1;
                            band_idx <= band_idx + 1; 
                        end
                        else begin
                            counter <= counter + 1;
                            eof     <= 0;
                        end
                    end
                    else begin
                        dout_valid <= 0;
                    end
                end
            endcase
        end
    end

    wire piso_ready;
    assign piso_ready = (state == READ);

    piso #(
        .DIN_WIDTH(DIN_WIDTH),
        .DOUT_WIDTH(512),
        .FIFO_DEPTH(FIFO_DEPTH)
    ) piso_inst (
        .clk(clk),
        .rst(rst),
        .din(din),
        .din_valid(din_valid),
        .dout(tge_data),
        .dout_valid(tge_data_valid),
        .dout_ready(piso_ready),
        .fifo_full(fifo_full)
    );

    assign tx_data  = dout;
    assign tx_valid = dout_valid;
    assign tx_eof   = eof;

endmodule
\end{lstlisting}

\begin{lstlisting}[style={prettyverilog}, caption={Verilog code for the top-level packetizer module used in the HGE write path.}, label={lst:verilog_packetizer_top}]
`include "/home/bruno/f-engine/simulink_models/includes.v"
module packetizer_top_tge
 (
    input wire clk,
    input wire ce,

    input wire rst,
    input wire [2047:0] din,
    input wire din_valid,
    
    //configuration signals
    input wire [31:0] pkt_len,
    input wire [31:0] sleep_cycles,
    input wire [31:0] timestamp_sec,  
    input wire [31:0] timestamp_usec, 


    //to the 100Gbe block
    output wire [511:0] tx_data,
    output wire tx_valid,
    output wire tx_eof,

    output wire fifo_full
);

tge_write_packetizer #(
    .DIN_WIDTH(2048),
    .FIFO_DEPTH(8192)
) tge_write_inst (
    .clk(clk),
    .rst(rst),
    .din(din),
    .din_valid(din_valid),
    .pkt_len(pkt_len),
    .sleep_cycles(sleep_cycles),
    .timestamp_sec(timestamp_sec),  
    .timestamp_usec(timestamp_usec), 
    .tx_data(tx_data),
    .tx_valid(tx_valid),
    .tx_eof(tx_eof),
    .fifo_full(fifo_full)
);

endmodule
\end{lstlisting}